\input{estilo}
% El mismo gráfico para cerrar la ventana de 8:00 a 8:05: cuando el watermark
% cruza el borde de las 8:05 (el punto rojo; en la pizarra, una cruz), se
% cierra, con los eventos atrapados entre la curva y la diagonal (verdes).
\begin{document}
\begin{figura}\begin{tikzpicture}
  \draw[->] (0,0) -- (7.2,0) node[etiq, below left=4pt and 0pt] {event time};
  \draw[->] (0,0) -- (0,5.8) node[etiq, above] {processing time};
  \draw[dashed] (0,0) -- (6.4,5.6);
  \draw[acento, line width=1.1pt] plot[smooth, tension=0.6] coordinates
    {(0,0) (0.6,1.5) (1.6,2.1) (2.2,2.4) (2.9,3.4) (3.4,4.1) (4.4,4.4) (5.4,5.0) (6.2,5.9)};
  \foreach \x/\t in {2.2/8:00, 3.4/8:05} {
    \draw[dashed, line width=0.5pt] (\x,0) -- (\x,4.9);
    \draw (\x,-0.12) -- (\x,0.12) node[etiq, below=6pt] {\t};
  }
  \draw[acento, line width=1.1pt] (2.2,4.1) -- (3.4,4.1);
  \draw[acento, line width=1.4pt] (3.4,-0.12) -- (3.4,0.12);
  \node[evento=acento, minimum size=6pt] at (3.4,4.1) {};
  \foreach \p in {(2.4,2.5),(2.65,2.95),(2.85,2.7),(3.05,3.4),(3.3,3.7)}
    \node[evento=segundo] at \p {};
\end{tikzpicture}\end{figura}
\end{document}
